
\textbf{Proposition 2.4.4.} \uline{L'application polynôme} @@M0@@ (2.3.1) \uline{est constante sur les classes suivant} @@M1@@.

\emph{Démonstration} : Si @@M2@@ est une extension algébriquement close de @@M3@@, on trouve un diagramme commutatif
@@M4@@
où @@M5@@ est l'application polynôme correspondant à @@M6@@ sur @@M7@@. On peut donc supposer @@M8@@ algébriquement clos. L'assertion découle alors immédiatement de 2.1.3 et 2.4.3.

\textbf{Remarque 2.4.5.} Dans Exp. XIII on va démontrer ce résultat sans utiliser la représentabilité du foncteur @@M9@@.

\section*{3. Formules de projection pour les gradués associés}

\subsection*{3.1.}
Considérons un schéma noethérien @@M10@@. On désigne par @@M11@@ le gradué associé à la @@M12@@-filtration de @@M13@@ et par @@M14@@ le gradué associé à la filtration 1.1.1 de @@M15@@. Grâce à 1.3.2, la loi de composition (1.3.1)
@@M16@@
induit par passage aux gradués associés des lois de composition
@@M17@@
